\documentclass[10pt,a4paper]{amsart}
\usepackage[margin=19mm]{geometry}
\usepackage{amsmath,amssymb,mathtools}
\usepackage[scr=rsfs,cal=euler]{mathalfa}
\usepackage{xfrac}
\usepackage[unicode,hidelinks,bookmarks=false]{hyperref}
\newcommand{\ie}{i.e.\ }
\newcommand{\cf}{cf.\ }
\newcommand{\resp}{respectively\ }
\newcommand{\Subsection}{\subsection}
\providecommand{\nobreakdash}{\nobreak\hbox{-}}
\setlength{\emergencystretch}{2em}
\multlinegap0pt
\usepackage{amssymb}
\usepackage{xcolor}
\usepackage{enumerate}
\newtheorem{proposition}{Proposition}
\newtheorem{pf}{Pf}
\title{Disjointly universal inner functions}
\author{\"Ozgur Martin, Dimitris Papathanasiou,, Sibel \c{S}ahin}
\date{Source: arXiv:2606.19225v1 (CC BY 4.0)}
\begin{document}
\maketitle

\noindent\emph{Source and licence.} Disjointly universal inner functions. Version arXiv:2606.19225v1 (CC BY 4.0), distributed by arXiv under the Creative Commons Attribution 4.0 International licence, http://creativecommons.org/licenses/by/4.0/, as recorded in the arXiv licence metadata for this exact version and shown on the abstract page https://arxiv.org/abs/2606.19225v1. Full licence notice: \texttt{licenses/arxiv-2606.19225-CC-BY-4.0.txt}.\par\smallskip

\noindent\textit{Editorial scope.} Counted unit: the article's proposition Birkhoff (source0.tex lines 248--255). The article's complete original proof of it is delivered with it (source0.tex lines 257--263, one \texttt{proof} environment of 7 lines). No mathematics is written, repaired or re-derived: the statement and the proof are the article's own lines, in the article's own notation, and no retained line is substituted or re-flowed.  No further source-local context is needed: every cross-reference of every retained line resolves inside the two delivered spans.  Nothing was omitted from any retained span, so the package carries no omission record.  No retained line contains a citation, so the wrapper carries no bibliography and no bibliography entry.  No retained line contains a figure environment, an image-inclusion command or drawing code, so no image is delivered and the package declares no asset dependency.  Editorial formatting only, in addition: an AMS \texttt{amsart} preamble that restates the article's own declarations and macros used by the retained lines, namely the environments \texttt{proposition}, \texttt{pf} on separate counters, together with the standard packages those lines need.  The retained environments print in this document as Proposition 1, Pf 1 in order of appearance, whereas the article prints the same items with the numbering of its own full document.  The complete native source of the article as distributed in the arXiv e-print for this exact version is bundled unchanged as \texttt{sources/203129/source0.tex}.
\begin{proposition} \label{Birkhoff}
Let $X$ be a separable, metrizable, complete topological space, 
%with identity $e$, 
and for each $n\in \mathbb{N}$, let $T_{1,n}$ and $T_{2,n}$ be continuous self maps of $X$. The set of disjointly universal elements for the sequences $(T_{1,n})_n, (T_{2,n})_n$  is a $G_{\delta}$ subset of $X$, and it is residual if and only if, for each $U, V_1$, and $V_2$ non-empty open subsets of $X$, there exists $N\in \mathbb{N}$ such that
\begin{equation}\label{Baire}
U \cap T^{-1}_{1,N} (V_1) \cap T^{-1}_{2,N} (V_2) \neq \emptyset.
\end{equation}
\end{proposition}

\begin{pf}
    Since $X$ is separable and metrizable, it is second countable. Let $(V_k)_{k=1}^{\infty}$ be a base for the topology of $X$. The set of disjointly universal vectors for $(T_{1,n})_n, (T_{2,n})_n$  can be written as
    $$
    \bigcap_{i,j=1}^{\infty}\bigcup_{n=1}^{\infty}( T^{-1}_{1,n} (V_i) \cap T^{-1}_{2,n} (V_j))
    $$
    which is a $G_{\delta}$ subset of $X$. Assuming equation \ref{Baire}, the Baire category theorem proves the residuality of the set of disjointly universal elements for $(T_{1,n})_n$, $(T_{2,n})_n$. The reverse implication is trivial.
    \end{pf}

\end{document}
